\chapter{Bandwidth, antibandwidth và tô đa màu}
\chapterlead{Các bài toán trong chương đều nói về khoảng cách giữa nhãn, nhưng
không chia sẻ một ngữ nghĩa: antibandwidth gán một hoán vị, tô màu theo
bandwidth gán một màu, còn tô đa màu gán một tập màu có thứ tự.}

\index{antibandwidth}\index{cyclic antibandwidth}
\index{bandwidth coloring}\index{bandwidth multicoloring}
\index{forbidden interval}

\flowdiagram{Hoán vị hay tô màu}{Cửa sổ/thứ tự/khối}
  {Khoảng cách}{Tìm cận}{Gán nhãn và chứng nhận}

\section{Bố trí đồ thị, gán nhãn và tô màu}

Từ ``bandwidth'' được dùng cho nhiều bài toán có công thức gần nhau nhưng lịch
sử thuật toán khác nhau. Bandwidth đồ thị cổ điển tối thiểu hóa khoảng cách
lớn nhất trên cạnh của một hoán vị đỉnh. Antibandwidth tối đa hóa khoảng cách
nhỏ nhất. Tô màu theo bandwidth bỏ tính song ánh và cho trọng số khoảng cách
trên cạnh; tô đa màu cho mỗi đỉnh nhiều tần số. Văn liệu gán nhãn đồ thị rộng
hơn còn chứa hàng trăm quy tắc gán số khác nhau
\parencite{gallian2024labeling}.

Ba truyền thống giải liên quan trực tiếp:
\begin{itemize}
  \item \textbf{cận lý thuyết đồ thị}: bậc, số ổn định, số sắc,
        \emph{clique} (đồ thị con đầy đủ) và cấu
        trúc lớp đồ thị;
  \item \textbf{phương pháp kinh nghiệm/siêu kinh nghiệm}: tìm kiếm cục bộ,
        \emph{tabu search} (tìm kiếm kiêng kỵ), lân cận biến đổi và thứ tự kiến tạo để tìm
        nghiệm đương nhiệm cho
        cận trên/cận dưới;
  \item \textbf{phương pháp chính xác}: quy hoạch động, nhánh--cận, MIP, CP và
        SAT theo cận.
\end{itemize}

SAT giải hiệu quả bài toán quyết định ``có phép gán nhãn tại cận \(k\) không''
và kết hợp tự nhiên với các cận lý thuyết đồ thị. Các cận thu hẹp khoảng tìm
kiếm, tạo các giả định đầu tiên tốt hơn và đôi khi trực tiếp chứng nhận tối ưu.
Với BCP/BMCP, khảo cứu tô màu đỉnh và các biến thể cung cấp mô hình, cận theo
đồ thị con đầy đủ
và nguồn gốc bộ kiểm chuẩn cần dùng làm đối chứng
\parencite{malaguti2010coloring}.

\begin{center}
\captionof{table}{Bốn họ bài toán khoảng cách nhãn.}
\label{tab:distance-families}
\begin{tabularx}{\textwidth}{@{}>{\bfseries\sffamily}p{.2\textwidth}YYY@{}}
\toprule
Họ & Phép gán & Cận miền tự nhiên & Đối chứng chính xác\\
\midrule
Bandwidth & Hoán vị & Cận bố trí/bậc & DP, B\&B, MIP\\
Antibandwidth & Hoán vị & Cận ổn định/số sắc & MIP, CP, SAT\\
BCP & Một màu/đỉnh & Cận đồ thị con đầy đủ và tô màu & MIP, CP, SAT\\
BMCP & Nhiều màu/đỉnh & Cận ô màu/nhu cầu & MIP, tìm kiếm cục bộ, SAT\\
\bottomrule
\end{tabularx}
\end{center}

\section{Phân loại trước khi biểu diễn}

\begin{center}
\captionof{table}{Phân biệt ngữ nghĩa của ABP, CABP, No-hole Anti-\(k\)-labeling, BCP và BMCP.}
\label{tab:distance-semantics}
\begin{tabularx}{\textwidth}{@{}>{\bfseries\sffamily}p{.24\textwidth}YYY@{}}
\toprule
Bài toán & Nghiệm & Ràng buộc & Mục tiêu\\
\midrule
Antibandwidth & Hoán vị \(1,\ldots,n\) &
\(\card{f(u)-f(v)}\ge k\) & Tối đa \(k\)\\
Antibandwidth tuần hoàn & Hoán vị trên vòng &
\(D_c(f(u),f(v))\ge k\) & Tối đa \(k\)\\
No-hole Anti-\(k\)-labeling & Ánh xạ toàn ánh \(\set{1,\ldots,k}\) (\(k<n\)) &
\(\card{\psi(u)-\psi(v)}\ge w\) và \(\sum_i x_{ij}\ge 1\) & Tối đa \(w\)\\
Tô màu theo bandwidth & Một màu mỗi đỉnh &
\(\card{c(u)-c(v)}\ge d_{uv}\) & Tối thiểu độ trải\\
Tô đa màu theo bandwidth & \(w(u)\) màu mỗi đỉnh &
Khoảng cách nội/ngoại đỉnh & Tối thiểu màu lớn nhất\\
\bottomrule
\end{tabularx}
\end{center}

Với ABP và CABP, nhãn là một hoán vị nên ma trận gán cần
\(\ExactlyOne\) theo cả hàng và cột. Với BCP, hai đỉnh không kề nhau có thể
dùng cùng màu; chỉ hàng cần \(\ExactlyOne\). Với BMCP, một đỉnh có nhiều ô
màu; thứ tự các ô là một phần của mô hình.

\begin{designrule}{Kiểm tra ngữ nghĩa trước AMO}
Hai phép biểu diễn có cùng biến \(x_{v,\ell}\) vẫn có thể mô tả hai bài toán khác
nhau. Hãy viết phép giải mã và ràng buộc đếm theo hàng/cột trước khi thêm các
ràng buộc khoảng cách.
\end{designrule}

\section{Antibandwidth và cửa sổ staircase}

\subsection{Cận và phương pháp chính xác ngoài SAT}

Với đồ thị liên thông, các cận trên ABP có thể lấy từ số đỉnh/cạnh, bậc, số ổn
định \(\alpha(G)\) và số sắc \(\chi(G)\). Chẳng hạn
\[
  AB(G)\le \alpha(G),
  \qquad
  AB(G)\le
  \left\lfloor\frac{n-1}{\chi(G)-1}\right\rfloor.
\]
Tính chính xác \(\alpha,\chi\) cũng khó, nhưng một cận trên cho
\(\alpha\) hoặc cận dưới cho \(\chi\) vẫn có thể tạo cận hợp lệ. Các mô hình
MIP, nhánh--cắt và CP đã cải thiện đáng kể chứng nhận trên bộ Harwell--Boeing
\parencite{sinnl2021antibandwidth}. Do đó một nghiên cứu SAT trên cùng bộ kiểm
chuẩn cần báo cùng cận đầu vào và so cả giá trị nghiệm lẫn khoảng cách chứng
minh, không chỉ số bài toán có mô hình.

\subsection[Tính khả thi ở cận k]{Tính khả thi ở cận \(k\)}

Cho đồ thị \(G=(V,E)\), \(\card{V}=n\). Biến
\[
  x_{v,\ell}=1\iff f(v)=\ell
\]
mô tả một hoán vị:
\[
  \ExactlyOne(x_{v,1},\ldots,x_{v,n})
  \quad\forall v,
\]
\[
  \ExactlyOne(x_{1,\ell},\ldots,x_{n,\ell})
  \quad\forall\ell.
\]
Với một cận \(k\), mỗi cạnh \((u,v)\) cấm hai nhãn cách nhau dưới \(k\):
\[
  \neg x_{u,a}\lor\neg x_{v,b},
  \qquad \card{a-b}<k.
\]

Một cách nhìn khác cố định vị trí \(a\) của \(u\). Khi đó \(v\) không được
nằm trong cửa sổ
\[
  [a-k+1,a+k-1]\cap[1,n].
\]
Khi \(a\) trượt, các cửa sổ chồng lấn tạo một họ AMO bậc thang. Nếu biểu diễn
từng cửa sổ độc lập, phần lớn bộ đếm bị lặp.

\begin{figure}[tb]
\centering
\begin{tikzpicture}[satfig,x=8mm,y=7mm]
  \node[satlabel] at (2.7,2.3) {một phép gán nhãn cho \(P_4\)};
  \foreach \i/\x/\lab in {1/0/2,2/1.8/4,3/3.6/1,4/5.4/3}{
    \node[satstate] (v\i) at (\x,1) {\(v_{\i}\)};
    \node[below=2pt of v\i,satlabel,text=Blue] {\(f(v_{\i})=\lab\)};
  }
  \draw[Ink,thick] (v1)--(v2)--(v3)--(v4);
  \node[satlabel] at (9,2.3) {cặp nhãn bị cấm khi \(k=2\)};
  \foreach \a in {1,...,4}{
    \node[left,satlabel] at (7.8,{1.6-.55*\a}) {\(\a\)};
    \node[above,satlabel] at ({7.8+.55*\a},1.6) {\(\a\)};
    \foreach \b in {1,...,4}{
      \pgfmathtruncatemacro{\delta}{abs(\a-\b)}
      \ifnum\delta<2
        \draw[Gold,fill=PaleGold] ({7.8+.55*\b},{1.6-.55*\a})
          rectangle ++(.5,-.5);
      \else
        \draw[Rule,fill=white] ({7.8+.55*\b},{1.6-.55*\a})
          rectangle ++(.5,-.5);
      \fi
    }
  }
  \node[satlabel,text=Gold] at (9,-1.65)
    {ô tô màu: \(|a-b|<2\)};
\end{tikzpicture}
\caption{Ở trái, nhãn \((2,4,1,3)\) cho \(P_4\) đạt antibandwidth \(2\).
Ở phải, mỗi ô được tô biểu diễn một cặp nhãn bị cấm trên một cạnh.}
\label{fig:abp-path-conflicts}
\end{figure}

\subsection{Thang Sequential Counter}

SCL chia dãy biến thành các khối, xây Sequential Counter trong mỗi khối rồi nối
trạng thái ở biên. Ý tưởng cốt lõi là biến phụ mang nghĩa số đếm tiền tố:
\[
  s_{i,j}=1\iff
  \sum_{q=1}^{i}x_q\ge j.
\]
Các cửa sổ chia sẻ tiền tố/hậu tố thay vì dựng lại bộ đếm. Với trường hợp
\(k=1\), đây là phép biểu diễn cho AMO bậc thang; với \AMK, bộ đếm giữ nhiều mức
\(j=1,\ldots,k\).

\begin{proposition}[Tái sử dụng qua biên khối]
Nếu bộ đếm trong mỗi khối đúng với mọi tiền tố cục bộ và bộ nối bảo toàn tổng
tích lũy khi đi qua biên, thì ngưỡng của mọi cửa sổ ghép từ hậu tố khối trái
và tiền tố khối phải đúng với số literal đúng trong cửa sổ đó.
\end{proposition}

\begin{proof}
Tổng của cửa sổ bằng tổng trên hai phần rời nhau. Bộ đếm cục bộ biểu diễn đúng
hai tổng; bộ nối liệt kê các cặp ngưỡng có tổng vượt cận. Vì vậy một phép gán
vi phạm ràng buộc đếm kích hoạt ít nhất một mệnh đề nối, còn một phép gán hợp
lệ mở rộng được bằng các giá trị đếm tiền tố thật.
\end{proof}

\textcite{truong2025scamo} cho thấy SCL giảm cả biến phụ và mệnh đề so với
nhiều phép biểu diễn không chia sẻ trên họ AMO bậc thang. Hướng \AMK được phát triển trong
\textcite{truong2026ladderamk}; Duplex của \textcite{fazekas2020duplex} là một
đối chiếu quan trọng theo hướng BDD hai chiều. Một cách nhìn rộng hơn về biến
phụ có cấu trúc được thảo luận trong \textcite{haberlandt2023aux}.

SCL bổ sung cho các mô hình MIP/CP một cơ chế nén các khoảng xung đột trong
\CNF. Cận mạnh và SCL tương hỗ trực tiếp: \(k\) quyết định chiều rộng cửa sổ
cần biểu diễn, còn Shared Counter giảm phần biểu diễn lặp giữa các cửa sổ đó.

\section{Antibandwidth tuần hoàn}

\subsection{Khoảng cách vòng}

Trong CABP, nhãn nằm trên vòng \(C_n\):
\[
  D_c(a,b)=\min\{\card{a-b},\,n-\card{a-b}\}.
\]
Với cận \(k\), một cạnh cấm cửa sổ vòng bán kính \(k-1\). Cửa sổ có thể vắt
qua điểm cắt nên dùng chỉ số môđun:
\[
  W(a,k)=\set{a-k+1,\ldots,a+k-1}\pmod n.
\]
Xung đột theo nhãn chính xác vẫn đúng nhưng số mệnh đề lớn. Cửa sổ vòng tạo
AMO tuần hoàn, có thể biểu diễn bằng cách cắt vòng tại một vị trí neo rồi nối hai
đầu bằng bộ nối.

\begin{figure}[tb]
\centering
\begin{tikzpicture}[satfig]
  \foreach \a/\ang in {1/90,2/45,3/0,4/-45,5/-90,6/-135,7/180,8/135}{
    \ifnum\a<4
      \node[satstate,satwarn] (c\a) at (\ang:24mm) {\(\a\)};
    \else
      \ifnum\a>6
        \node[satstate,satwarn] (c\a) at (\ang:24mm) {\(\a\)};
      \else
        \node[satstate,sattrue] (c\a) at (\ang:24mm) {\(\a\)};
      \fi
    \fi
  }
  \draw[Rule,thick] (c1)--(c2)--(c3)--(c4)--(c5)--(c6)--(c7)--(c8)--cycle;
  \node[satlabel,align=center] at (0,0)
    {\(W(1,3)=\{7,8,1,2,3\}\)\\vắt qua điểm cắt};
  \draw[Gold,very thick,->] (c1) to[bend left=20] (c3);
  \draw[Gold,very thick,->] (c1) to[bend right=20] (c7);
\end{tikzpicture}
\caption{Cửa sổ cấm trên vòng \(C_8\) với tâm \(1\) và \(k=3\). Hai nhánh
\(7,8\) và \(2,3\) phải được nối khi tuyến tính hóa vòng.}
\label{fig:cabp-cyclic-window}
\end{figure}

\subsection{Đối xứng của vòng}

CABP bất biến dưới phép quay và phản xạ của nhãn. Ta có thể cố định một đỉnh
đại diện tại nhãn \(1\), rồi chọn hướng chuẩn cho một đỉnh thứ hai. Đây là phá
đối xứng ở cấp hoán vị; điều kiện này độc lập với SCL dùng cho các cửa sổ khoảng cách.

Tìm nghiệm tối ưu theo cận tăng:
\[
  F(k)=\SAT \Longrightarrow F(k-1)=\SAT,
  \qquad
  F(k)=\UNSAT \Longrightarrow F(k+1)=\UNSAT.
\]
Một nghiệm tại \(k^\star\) và \UNSAT tại \(k^\star+1\) chứng nhận tối ưu.

\begin{resultbox}{CABP và bộ kiểm chuẩn}
Nghiên cứu của \textcite{truong2026cabp} xây dựng các phép biểu diễn SAT cho CABP
và khai thác cửa sổ tuần hoàn cùng đối xứng của vòng. Một tập công thức khả thi được
chuẩn bị cho SAT Competition 2026 \parencite{truong2026competition}; về mặt
thiết kế, giá trị của bộ này là tách bài toán tối ưu thành chuỗi bài toán
\SAT/\UNSAT có cận \(k\) rõ ràng.
\end{resultbox}

\section{Bài toán No-hole Anti-\(k\)-labeling và nén khoảng cấm bằng biến tiền tố}

\index{No-hole Anti-k-labeling}\index{anti-k-labeling}
\index{prefix variables}\index{forbidden interval}

\subsection{Bài toán gán nhãn phi song ánh và điều kiện toàn ánh}

Trong quản lý phổ tần số và tài nguyên mạng không dây, số lượng kênh truyền dẫn khả thi \(k\) thường nhỏ hơn đáng kể so với số lượng trạm phát sóng \(n = \card{V}\) (\(k < n\)) do các giới hạn về mặt pháp lý và kỹ thuật. Bối cảnh này dẫn tới bài toán \emph{Anti-\(k\)-labeling} \parencite{guan2019anti}, một mở rộng phi song ánh tự nhiên của bài toán Antibandwidth cổ điển (nơi \(k=n\)). Khi \(k < n\), phép gán nhãn không còn là một hoán vị mà trở thành một phép phân hoạch các đỉnh vào các lớp nhãn.

Biến thể \emph{No-hole Anti-\(k\)-labeling} \parencite{wang2021anti,hai2026nohole} yêu cầu hàm gán nhãn \(\psi: V \to \set{1,\ldots,k}\) phải là một phép ánh xạ \emph{toàn ánh} (\emph{surjective mapping}), tức là mọi nhãn trong tập \(\set{1,\ldots,k}\) đều phải được gán cho ít nhất một đỉnh (không có nhãn trống/lỗ hổng). Mục tiêu của bài toán là tìm hàm gán nhãn \(\psi\) sao cho khoảng cách giữa hai đỉnh kề nhau là lớn nhất:
\begin{equation}
  \lambda_k^{nh}(G) = \max_{\psi} \min_{(u,v)\in E} \card{\psi(u) - \psi(v)}.
  \label{eq:nohole-objective}
\end{equation}

Mô hình quy hoạch tuyến tính số nguyên (ILP) cho bài toán quyết định ``liệu có tồn tại phép gán nhãn toàn ánh thỏa mãn khoảng cách tối thiểu \(w\) không'' sử dụng các biến nhị phân \(x_{ij} = 1 \iff \psi(i) = j\):
\begin{align}
  \sum_{j=1}^k x_{ij} &= 1 && \forall i\in\set{1,\ldots,n}, \label{eq:nohole-ilp-vertex}\\
  \sum_{i=1}^n x_{ij} &\ge 1 && \forall j\in\set{1,\ldots,k}, \label{eq:nohole-ilp-nohole}\\
  x_{uj} + \sum_{j'=\max(1, j-w+1)}^{\min(k, j+w-1)} x_{vj'} &\le 1 && \forall (u,v)\in E, j\in\set{1,\ldots,k}. \label{eq:nohole-ilp-distance}
\end{align}
Ràng buộc~\eqref{eq:nohole-ilp-vertex} bảo đảm mỗi đỉnh nhận đúng một nhãn; Ràng buộc~\eqref{eq:nohole-ilp-nohole} bảo đảm tính toàn ánh (\emph{no-hole}); và Ràng buộc~\eqref{eq:nohole-ilp-distance} cấm hai đỉnh kề nhau nhận các nhãn thuộc khoảng cấm \((j-w, j+w)\).

\subsection{Từ bùng nổ mệnh đề từng cặp đến biểu diễn khoảng bằng biến tiền tố}

Chuyển trực tiếp Ràng buộc khoảng cách~\eqref{eq:nohole-ilp-distance} sang \CNF bằng phép biểu diễn từng cặp (Direct Encoding - DE) sinh ra họ mệnh đề xung đột:
\begin{equation}
  \neg x_{ui} \lor \neg x_{vj} \qquad \forall (u,v)\in E, i\in\set{1,\dots,k}, j\in [\max(1,i-w+1),\min(k,i-w-1)].
  \label{eq:nohole-direct-cnf}
\end{equation}
Số lượng mệnh đề của DE phát triển theo cấp số \(O(\card{E}\cdot k\cdot w)\). Khi khoảng cách mục tiêu \(w\) tăng lên, số mệnh đề khoảng cách bùng nổ mạnh mẽ, gây tắc nghẽn bộ nhớ và làm suy giảm hiệu năng của bộ giải SAT.

Để giải quyết triệt để nút thắt này, phép biểu diễn \mbox{\textnormal{\textsc{PSE}}} (\emph{Proposed SAT Encoding}) \parencite{hai2026nohole} giới thiệu các biến phụ tiền tố \(l_{ij} \in \set{0,1}\) với ngữ nghĩa:
\[
  l_{ij} = 1 \iff \psi(i) \le j \iff \exists k' \le j : x_{ik'} = 1.
\]
Tính đơn điệu tự nhiên của biến tiền tố đòi hỏi \(l_{ij} \to l_{i(j+1)}\) cho mọi \(j < k\). Hệ thống mệnh đề liên kết biến gán gốc \(x_{ij}\) và biến tiền tố \(l_{ij}\) được định nghĩa bởi:
\begin{align}
  x_{ij} &\to l_{ij} && \forall i\in\set{1,\ldots,n}, j\in\set{1,\ldots,k}, \label{eq:nohole-link-1}\\
  l_{ij} &\to l_{i(j+1)} && \forall i\in\set{1,\ldots,n}, j\in\set{1,\ldots,k-1}, \label{eq:nohole-link-2}\\
  \neg x_{i1} &\to \neg l_{i1} && \forall i\in\set{1,\ldots,n}, \label{eq:nohole-link-3}\\
  (\neg x_{ij} \land \neg l_{i(j-1)}) &\to \neg l_{ij} && \forall i\in\set{1,\ldots,n}, j\in\set{2,\ldots,k}, \label{eq:nohole-link-4}\\
  x_{ij} &\to \neg l_{i(j-1)} && \forall i\in\set{1,\ldots,n}, j\in\set{2,\ldots,k}. \label{eq:nohole-link-5}
\end{align}

Nhờ cấu trúc biến tiền tố, toàn bộ khoảng cấm \((j-w, j+w)\) giữa hai đỉnh kề nhau \((u,v)\in E\) được nén lại chỉ trong **đúng một mệnh đề kéo theo đơn duy nhất**:
\begin{equation}
  x_{uj} \to \left( l_{v, \max(1, j-w)} \lor \neg l_{v, \min(k, j+w-1)} \right) \qquad \forall (u,v)\in E, j\in\set{1,\ldots,k}.
  \label{eq:nohole-prefix-distance}
\end{equation}
Mệnh đề~\eqref{eq:nohole-prefix-distance} khẳng định rằng: nếu đỉnh \(u\) nhận nhãn \(j\), thì đỉnh kề \(v\) bắt buộc phải nhận nhãn nhỏ hơn/bằng \(j-w\) (thỏa mãn literal \(l_{v, \max(1,j-w)}\)) hoặc phải nhận nhãn lớn hơn \(j+w-1\) (thỏa mãn literal \(\neg l_{v, \min(k,j+w-1)}\)). Điều này tự động loại bỏ tất cả các nhãn nằm trong khoảng cấm \([j-w+1, j+w-1]\) mà không cần liệt kê từng cặp nhãn vi phạm.

\begin{figure}[tb]
\centering
\begin{tikzpicture}[satfig,x=8.5mm,y=7.5mm]
  %--- Direct Encoding ---
  \node[anchor=west,font=\rmfamily\bfseries\small] at (0,3.5) {(a) Phép biểu diễn từng cặp (Direct Encoding)};
  \node[anchor=east,font=\rmfamily\small] at (0.6,2.1) {$x_{u,3}=1 \implies$};
  \foreach \j/\v/\c in {1/1/PaleTeal, 2/2/PaleGold, 3/3/PaleGold, 4/4/PaleGold, 5/5/PaleTeal} {
    \draw[satbox,fill=\c] (0.8+\j*1.1, 1.75) rectangle +(0.8,0.7);
    \node[font=\rmfamily\small] at (0.8+\j*1.1+0.4, 2.1) {$\v$};
  }
  \node[anchor=west,font=\rmfamily\small,text=Gold!80!black] at (0.8, 0.8) {Liệt kê 3 mệnh đề cấm độc lập: $\neg x_{v,2}$, $\neg x_{v,3}$, $\neg x_{v,4}$};

  %--- Separator ---
  \draw[Rule,thick] (7.2,0.1)--(7.2,3.6);

  %--- Proposed Prefix Encoding ---
  \node[anchor=west,font=\rmfamily\bfseries\small] at (7.8,3.5) {(b) Phép biểu diễn tiền tố khoảng cấm (PSE)};
  \node[anchor=east,font=\rmfamily\small] at (8.4,2.1) {$x_{u,3}=1 \implies$};
  
  \draw[satbox,fill=PaleTeal,draw=Teal,thick] (8.6,1.75) rectangle +(0.8,0.7);
  \node[font=\rmfamily\small] at (9.0,2.1) {$1$};
  \node[anchor=south,font=\rmfamily\small,text=Teal] at (9.0,2.55) {$l_{v,1}=1$};

  \draw[satbox,fill=SoftGray,dashed] (9.6,1.75) rectangle +(3.0,0.7);
  \node[font=\rmfamily\small,text=black!60] at (11.1,2.1) {Khoảng cấm $[2, 4]$};

  \draw[satbox,fill=PaleTeal,draw=Teal,thick] (12.8,1.75) rectangle +(0.8,0.7);
  \node[font=\rmfamily\small] at (13.2,2.1) {$5$};
  \node[anchor=south,font=\rmfamily\small,text=Teal] at (13.2,2.55) {$\neg l_{v,4}=1$};

  \node[anchor=west,font=\rmfamily\small,text=Teal] at (8.6,0.8) {Một mệnh đề đơn duy nhất: $x_{u,3} \to (l_{v,1} \lor \neg l_{v,4})$};
\end{tikzpicture}
\caption{So sánh cơ chế nén khoảng cấm khi $k=5, w=2$: (a) Phép biểu diễn từng cặp liệt kê $O(w)$ mệnh đề; (b) Phép biểu diễn tiền tố PSE loại toàn bộ khoảng cấm bằng một mệnh đề kéo theo duy nhất.}
\label{fig:nohole-prefix-interval}
\end{figure}

\begin{proposition}[Độ phức tạp và tính đúng đắn của PSE]
\label{prop:nohole-pse-complexity}
Phép biểu diễn \mbox{\textnormal{\textsc{PSE}}} biểu diễn chính xác bài toán No-hole Anti-\(k\)-labeling. Tổng số biến nhị phân là \(2nk\), tổng số mệnh đề khoảng cách là \(2\card{E}k\), và các mệnh đề liên kết là \(n(4k-2)\). Phép biểu diễn hoàn toàn triệt tiêu sự phụ thuộc của số mệnh đề vào tham số khoảng cách mục tiêu \(w\).
\end{proposition}

\begin{proof}
Mỗi đỉnh \(i\) có \(k\) biến gán gốc \(x_{ij}\) và \(k\) biến tiền tố \(l_{ij}\), tổng cộng \(2nk\) biến. Các Mệnh đề~\eqref{eq:nohole-link-1}--\eqref{eq:nohole-link-5} bảo đảm ngữ nghĩa \(l_{ij} = \sum_{k'\le j}x_{ik'}\) và ép buộc đỉnh \(i\) nhận đúng một nhãn mà không cần thêm các mệnh đề \textsc{ExactlyOne} bổ sung. Mệnh đề khoảng cách~\eqref{eq:nohole-prefix-distance} chỉ tạo ra 2 mệnh đề cho mỗi cặp cạnh định hướng và nhãn \(j\), cho tổng số \(2\card{E}k\) mệnh đề.
\end{proof}

\subsection{Phá vỡ đối xứng phản xạ và kết quả thực nghiệm}

Không gian nghiệm của bài toán No-hole Anti-\(k\)-labeling chứa đối xứng phản xạ: nếu \(\psi(v)\) là một nghiệm hợp lệ, thì nghiệm đối xứng \(\psi'(v) = k + 1 - \psi(v)\) cũng hợp lệ và giữ nguyên khoảng cách \(\card{\psi'(u) - \psi'(v)} = \card{\psi(u) - \psi(v)}\) \parencite{hai2026nohole}. Để phá vỡ đối xứng này và thu hẹp không gian tìm kiếm, phép biểu diễn cố định một đỉnh chốt \(p\) (thường chọn đỉnh có chỉ số 0) thuộc nửa dưới miền nhãn:
\begin{equation}
  \bigwedge_{i=\lceil k/2 \rceil + 1}^k \neg x_{pi}.
  \label{eq:nohole-symmetry-breaking}
\end{equation}

\begin{resultbox}{Kết quả tiêu biểu cho No-hole Anti-\(k\)-labeling}
Trên 23 bộ dữ liệu chuẩn Harwell--Boeing \parencite{hai2026nohole}:
\begin{itemize}
  \item Phép biểu diễn \mbox{\textnormal{\textsc{PSE}}} kết hợp với bộ giải song song PAINLESS giải thành công \textbf{100\% các bài toán}, vượt trội hoàn toàn so với mô hình \mbox{\textnormal{\textsc{DE}}} (bị bùng nổ mệnh đề và không tìm được nghiệm khả thi trên 3 bài toán lớn U, V, X trong thời gian 1800 giây).
  \item So với các công cụ thương mại hàng đầu đại diện cho MIP (Gurobi, CPLEX) và CP (CPLEX CP)---vốn hoàn toàn thất bại không tìm được bất kỳ nghiệm khả thi nào trên 8 bài toán quy mô lớn từ M đến X---phép biểu diễn \mbox{\textnormal{\textsc{PSE}}} cải thiện trung bình \textbf{19\%} chất lượng nghiệm mục tiêu và giảm trên \textbf{82\%} thời gian chạy trên các bài toán giải được.
\end{itemize}
\end{resultbox}

\section{Tô màu theo bandwidth}

\subsection{Quan hệ với tô màu đỉnh}

Khi \(d_{uv}=1\) trên mọi cạnh, BCP trở về tô màu đỉnh theo độ trải. Khi các
trọng số lớn hơn, một đồ thị con đầy đủ không chỉ yêu cầu các màu khác nhau mà
còn yêu cầu một phép đóng gói một chiều của các khoảng cách. Vì vậy kích thước
đồ thị con đầy đủ
vẫn cho cận dưới nhưng không mô tả đủ phép tách có trọng số. Các mô hình MIP
theo phép gán, đại diện hoặc phủ tập trong tô màu đỉnh là điểm khởi đầu; BCP
phải bổ sung các xung đột có xét khoảng cách
\parencite{malaguti2010coloring}.

Từ góc nhìn biểu diễn, biến trực tiếp mở rộng mô hình gán; biến thứ tự nén một dãy
giá trị xung đột thành hai biên; biến khối tạo các đại diện khoảng. Sáu họ dưới
đây khảo sát có hệ thống trục biểu diễn đó và kết nối với các cận lý thuyết đồ
thị đã nêu.

\subsection{Mô hình và sáu họ biểu diễn}

BCP cho đồ thị có trọng số cạnh \(d_{uv}>0\). Mỗi đỉnh nhận màu
\(c(u)\in[1,H]\) và
\[
  \card{c(u)-c(v)}\ge d_{uv}
  \quad\forall uv\in E.
\]
Mục tiêu là tối thiểu
\[
  \chi_B(G)=\min_c\max_{u\in V}c(u).
\]

Nghiên cứu hệ thống của \textcite{nguyen2026bcp} tổ chức không gian thiết kế thành sáu
họ:
\begin{center}
\captionof{table}{Sáu họ biểu diễn cho tô màu theo bandwidth.}
\label{tab:bcp-encodings}
\begin{tabularx}{\textwidth}{@{}>{\bfseries\sffamily}p{.14\textwidth}p{.31\textwidth}Y@{}}
\toprule
Họ & Biến & Cơ chế khoảng cách\\
\midrule
1G & \(g_{u,j}\iff c(u)\ge j\) & Một chuỗi thứ tự hướng lớn hơn\\
1L & \(l_{u,j}\iff c(u)\le j\) & Một chuỗi thứ tự hướng nhỏ hơn\\
2G & Giá trị đúng \(x_{u,j}\) và ngưỡng \(g_{u,j}\) &
Giá trị chính xác kích hoạt hai biên khoảng cấm\\
2L & Giá trị đúng \(x_{u,j}\) và ngưỡng \(l_{u,j}\) &
Như 2G với hướng thứ tự đảo\\
X & Biến khối & Chia sẻ cấu trúc khoảng cách theo khối\\
\(X_a\) & Khối và phép trừ phụ &
Tái sử dụng hiệu/khoảng qua biến phụ\\
\bottomrule
\end{tabularx}
\end{center}

\subsection{Khoảng cấm bằng hai biên}

Giả sử \(x_{u,a}=1\). Khi đó màu của \(v\) phải nằm ngoài
\[
  [a-d_{uv}+1,\;a+d_{uv}-1].
\]
Với ngữ nghĩa \(g_{v,t}=1\iff c(v)\ge t\), điều kiện được viết bằng một mệnh đề
biên:
\[
  \neg x_{u,a}
  \lor
  \neg g_{v,a-d_{uv}+1}
  \lor
  g_{v,a+d_{uv}}.
\]
Hai literal cuối lần lượt nói \(c(v)<a-d_{uv}+1\) hoặc
\(c(v)\ge a+d_{uv}\). Số mệnh đề không còn phụ thuộc vào độ rộng khoảng cấm;
mỗi giá trị chính xác tạo một mệnh đề cho mỗi hàng xóm.

Biểu diễn theo khối chia miền màu thành các đoạn và dùng biến phụ để tóm tắt quan
hệ giữa các đoạn. CNF có thể lớn hơn Order Encoding, nhưng lan truyền trên bài
toán khó có thể mạnh hơn vì bộ giải nhìn thấy các cấu trúc khoảng được chia sẻ.

\begin{workedexample}{Một khoảng cấm được nén thành hai biên}
Cho \(H=7\), \(x_{u,4}=1\) và \(d_{uv}=2\). Đỉnh \(v\) không được nhận
\[
  [4-2+1,\,4+2-1]=[3,5].
\]
Với \(g_{v,t}\Leftrightarrow c(v)\ge t\), chỉ cần mệnh đề
\[
  \neg x_{u,4}\lor\neg g_{v,3}\lor g_{v,6}.
\]
Khi \(x_{u,4}=1\), hai literal còn lại nói đúng hai miền hợp lệ:
\(c(v)\le2\) hoặc \(c(v)\ge6\). Direct Encoding từng cặp phải cấm riêng
\((4,3),(4,4),(4,5)\); biểu diễn theo biên giữ số mệnh đề không đổi khi khoảng
cấm rộng hơn, miễn tập cấm vẫn là một khoảng.
\end{workedexample}

\begin{resultbox}{Hiệu ứng tương tác}
Trên tập GEOM và MS-CAP của \textcite{nguyen2026bcp}, cấu hình theo khối giải
được các bài toán khó mà Order Encoding không hoàn tất trong cùng giới hạn.
Giải tăng dần và đối xứng không có hiệu ứng đồng đều: lợi ích phụ thuộc vào họ
biểu diễn. Vì vậy thí nghiệm phân tích thành phần (ablation study) nên dùng đơn vị
\(\text{phép biểu diễn}\times\text{đặc trưng}\), không chỉ so trung bình từng
đặc trưng.
\end{resultbox}

\section{Tô đa màu theo bandwidth}

\subsection{Mô hình ô màu và giảm miền}

BMCP gán cho đỉnh \(u\) một tập \(w(u)\) màu. Viết chúng thành các ô có thứ
tự
\[
  c(u)_1<c(u)_2<\cdots<c(u)_{w(u)}.
\]
Hai ô cùng đỉnh cách nhau ít nhất \(d(u,u)\); hai ô ở hai đỉnh kề nhau cách ít
nhất \(d(u,v)\). Với cận màu lớn nhất \(k\), miền ô được siết:
\[
  LB_{u,i}=1+(i-1)d(u,u),
\]
\[
  UB_{u,i}=k-(w(u)-i)d(u,u).
\]
Thứ tự ô loại đối xứng hoán vị màu bên trong một đỉnh và đồng thời tạo lan
truyền cận.

Mô hình ô là một phép mở rộng miền, trong khi mô hình tần số là một mô hình
tập. Hai cách có đối xứng và cận dưới khác nhau. Trong mô hình ô, thứ tự nội bộ
loại \(w(u)!\) hoán vị; trong mô hình tập, đối xứng đó không được tạo ra nhưng
ràng buộc ``chọn đúng \(w(u)\) màu'' cần ràng buộc đếm. So sánh FE/SBE/OBE vì
vậy phải bao gồm cả số trạng thái gán trước các mệnh đề khoảng cách.

\subsection{FE, SBE và OBE}

\paragraph{Flat Encoding (biểu diễn phẳng, FE).}
Biến \(x_{u,m}\) cho biết màu \(m\) được dùng tại \(u\). Ràng buộc đếm bảo đảm
đúng \(w(u)\) màu; xung đột từng cặp cấm các cặp quá gần.

\paragraph{Slot-Based Encoding (biểu diễn theo ô màu, SBE).}
Biến \(x_{u,i,m}\) nói ô \(i\) bằng \(m\). Mỗi ô có ràng buộc đúng một và các
cặp ô vi phạm khoảng cách bị cấm trực tiếp.

\paragraph{Order-Based Encoding (biểu diễn theo thứ tự, OBE).}
Thêm \(g_{u,i,m}=1\iff c(u)_i\ge m\). Giá trị chính xác tại một ô kích hoạt
hai biên khoảng cấm của ô ở đỉnh khác:
\[
  x_{u,i,m}\rightarrow
  \left(
    g_{v,j,m+d(u,v)}
    \lor
    \neg g_{v,j,m-d(u,v)+1}
  \right).
\]
Số mệnh đề xung đột giảm một bậc theo miền màu so với SBE từng cặp.

\begin{proposition}[Đúng đắn của OBE]
Nếu liên kết giá trị chính xác--thứ tự xác định cùng một giá trị cho mỗi ô,
chuỗi thứ tự đơn điệu và mọi giá trị chính xác kích hoạt mệnh đề khoảng cấm ở
trên, thì mọi mô hình
OBE giải mã thành một phép tô đa màu theo bandwidth hợp lệ.
\end{proposition}

\begin{proof}
Ràng buộc đúng một cho mỗi ô một màu; các mệnh đề cận và thứ tự giữ các ô trong
miền và đúng thứ tự. Với mỗi cặp ô bị ràng buộc, giá trị chính xác \(m\) của ô
thứ nhất buộc ô thứ hai nằm dưới \(m-d+1\) hoặc từ \(m+d\) trở lên. Vì vậy
khoảng cách tuyệt đối ít nhất \(d\).
\end{proof}

Tìm kiếm tăng dần từ trên xuống bắt đầu từ một cận trên khả thi. Sau khi tìm
mô hình có độ trải \(s\), bộ giải cấm các màu từ \(s\) trở lên và tiếp tục.
Lần \UNSAT đầu tiên tại \(s-1\), cùng mô hình tại \(s\), chứng nhận tối ưu.
\textcite{nguyen2026bmcp} dùng cơ chế này để chứng minh tối ưu cho BMCP.

\section{Một trừu tượng chung: giá trị chính xác và khoảng cấm}

ABP, BCP và BMCP gặp cùng một mẫu:
\begin{enumerate}
  \item một biến chính xác chọn giá trị \(a\);
  \item lựa chọn đó tạo một khoảng giá trị bị cấm cho đối tượng khác;
  \item biến thứ tự nén khoảng thành hai literal biên.
\end{enumerate}
Khác biệt nằm ở nguồn của khoảng và số lựa chọn trong phép gán. ABP có ràng
buộc khác nhau toàn cục; BCP không có; BMCP có nhiều ô trên một đỉnh. Vì vậy
mô-đun \texttt{ForbiddenInterval} có thể dùng lại, còn
\texttt{AssignmentSemantics} phải được tham số hóa theo bài toán.

\section{Quy trình chính xác kết hợp}

Một quy trình có giá trị tham khảo cho cả bốn bài toán là:
\begin{enumerate}
  \item dùng cận lý thuyết đồ thị và phương pháp kinh nghiệm để tạo khoảng
        \([LB,UB]\);
  \item chọn mô hình đúng ngữ nghĩa: hoán vị, một màu hoặc nhiều ô;
  \item lọc miền bằng cận cục bộ hoặc đồ thị con đầy đủ trước khi tạo biến;
  \item chọn Direct Encoding/thứ tự/khối theo hình dạng tập cấm;
  \item giải theo quy trình cận đơn điệu và lưu nghiệm làm chứng;
  \item tại cận cuối, kiểm tra phép gán nhãn miền và chứng minh \UNSAT độc lập.
\end{enumerate}

Kiến trúc này đặt SAT trong một khung giải chính xác thay vì coi bộ giải SAT là
toàn bộ thuật toán. Cấu trúc này cũng làm rõ vị trí đóng góp của một nghiên cứu mới: cận, mô
hình, phép biểu diễn, quy trình tìm kiếm hay chứng nhận.

\begin{summarybox}
\textbf{Bài học trung tâm.} Chọn Direct Encoding, thứ tự hay theo khối không
chỉ là chọn kích thước \CNF. Biến trực tiếp làm phép giải mã và xung đột cục bộ
rõ; biến thứ tự nén bất đẳng thức và khoảng; biến khối tạo điểm chia sẻ mới cho
lan truyền. Một kiến trúc tốt thường kết hợp cả ba ở những lớp khác nhau.
\end{summarybox}

\subsection*{Câu hỏi nghiên cứu}
\begin{enumerate}
  \item Với ABP trên đường \(P_4\), dựng \(F(k)\) tại \(k=2\) và chỉ rõ
        exactly-one theo hàng lẫn cột.
  \item Dẫn xuất mệnh đề biên BCP cho ba trường hợp: khoảng nằm trong
        miền, chạm biên trái và chạm biên phải.
  \item Thiết kế thí nghiệm phân tích thành phần (ablation study) \(2\times2\) cho giải tăng dần và đối xứng
        trên hai họ biểu diễn BCP.
\end{enumerate}
